\subsection{Relative norm repair with the prior radical fields retained}

\begin{statement}{A.117}{Proved}
{Finite approximation with prescribed splitting and a free coefficient place}
Retain an actual finite stage of A.44 and write its coefficients as
\[
 a=3u^{24},\qquad b=-2v^{24},\qquad u,v\in\Q^\times,
 \qquad v_2(u)=v_3(u)=v_2(v)=v_3(v)=0.
 \tag{A.117.1}
\]
Let \(n\ge1\), \(N=3^n\), \(d=N/3\), and choose
\(\alpha^N=a\).  Put \(K_n=\Q(\alpha)\),
\(K_{n-1}=\Q(\alpha^3)\), with \(K_0=\Q\).
Use an actual representative \(\beta\in K_n^\times\) of rational
norm \(b\) with the canonical two-adic divisor of A.101.
For integers \(h,r\ge1\), \(t\ge n\), and \(M\) divisible by
\(\operatorname{lcm}(24,N)\), choose
\[
 \begin{gathered}
 R=\Q(\mu_D),\qquad
       \operatorname{lcm}(M,3^h,3^r,3^t)\mid D,\\
 \theta^{3^t}=a,\quad \alpha=\theta^{3^{t-n}},
 \quad \beta_i=\tau^i\beta,\quad
       \tau(\alpha)=\zeta_N\alpha,\\
 E=R\bigl(\theta,\beta_i^{1/3^h}\ (i\bmod N),b^{1/M}\bigr).
 \end{gathered}
 \tag{A.117.2}
\]
These are the actual simultaneous fields of A.101 and A.105.
Define the relative torus and its finite modules by
\[
 \begin{gathered}
 T_n=\ker\bigl(\operatorname{Res}_{K_n/\Q}\mathbf G_m
          \xrightarrow{N_{K_n/K_{n-1}}}
                         \operatorname{Res}_{K_{n-1}/\Q}\mathbf G_m\bigr),\\
 R_r=\Z/3^r\Z,\qquad P_j=R_r[\Z/j\Z],\\
 C_r=T_n[3^r]=\ker(P_N\longrightarrow P_d)\otimes\mu_{3^r},
 \qquad C_r^\vee=P_N/\iota P_d .
 \end{gathered}
 \tag{A.117.3}
\]
The first permutation map sums the three entries in each block;
\(\iota\) repeats each entry on that same block.
For a decomposition group \(D_3\subset\operatorname{Gal}(E/\Q)\),
restriction is injective for both \(V=C_r\) and \(V=C_r^\vee\):
\[
 H^1(\operatorname{Gal}(E/\Q),V)
                 \longrightarrow H^1(D_3,V).
 \tag{A.117.4}
\]

Let \(S\) be any finite set containing the real place, three,
and every prime ramified in \(E\).  Let \(\mathcal P\) be the
primes outside \(S\) splitting completely in \(E\).
Write \(H^1_{S\cup\mathcal P}(\Q,V)\) for the union of
\(H^1(G_{\Q,S\cup P},V)\) over finite \(P\subset\mathcal P\).
Then
\[
 H^1_{S\cup\mathcal P}(\Q,V)\longrightarrow
                  \prod_{w\in S\setminus\{3\}}H^1(\Q_w,V)
             \quad\hbox{is onto}.
 \tag{A.117.5}
\]
A finite set \(P\) already suffices for this surjection.
Its primes may avoid any fixed finite set and lie above any
fixed bound.

For \(V=C_r\), every prescription in (A.117.5) consisting of
local Kummer classes of actual relative norm-one elements is
realized by an actual
\[
 U\in K_n^\times,\qquad N_{K_n/K_{n-1}}U=1.
 \tag{A.117.6}
\]
Its class \(\kappa_{T_n,r}(U)\) is unramified outside
\(S\cup P\) and has the prescribed classes at \(S\setminus\{3\}\).
Thus (A.117.6) preserves the lower norm exactly.
The assertion concerns the ramification of this finite Kummer
class; it does not assert that the rational divisor of a chosen
representative is supported on \(S\cup P\).
The three-adic class is free in (A.117.5), and no cyclic-image
condition is asserted at the newly permitted primes.
\end{statement}

\begin{proof}
\emph{The exact finite modules and the deeper splitting group.}
The torsion of the coordinatewise relative norm is the kernel
in (A.117.3).  The coordinate pairing identifies its Cartier
dual with the displayed quotient, since the annihilator of
the block-sum kernel is exactly the block-constant submodule.
Set
\[
 F=R(\theta),\qquad
 G_0=\operatorname{Gal}(F/\Q)
       =C_{3^t}\rtimes\operatorname{Gal}(R/\Q).
 \tag{A.117.7}
\]
The full Kummer degree and the literal splitting fixing
\(\theta\) were proved in A.105.  The translation acts on
the two modules through \(x\), of order \(N\).
Put \(s(x)=1+x^d+x^{2d}\).
As translation modules both are
\[
                         R_r[x]/s(x).
 \tag{A.117.8}
\]
For \(C_r^\vee\) this is the quotient by the block-constant
ideal \((s)\).  For \(C_r\), multiplication by \(x^d-1\)
identifies this module with the block-sum kernel in
\(R_r[x]/(x^N-1)\), followed by the cyclotomic twist.
The identification is injective because the monic polynomial
\(x^d-1\) can be cancelled from
\((x^d-1)s=x^N-1\) in polynomial division over \(R_r\).

Let \(\epsilon\) be complex conjugation in
\(\operatorname{Gal}(R/\Q)\).  On the presentation (A.117.8)
its two actions are, literally,
\[
 \epsilon f(x)=f(x^{-1})\quad(C_r^\vee),\qquad
 \epsilon f(x)=x^{-d}f(x^{-1})\quad(C_r).
 \tag{A.117.9}
\]
The second formula uses both
\((x^{-d}-1)/(x^d-1)=-x^{-d}\) and the cyclotomic sign.
In particular the involution acts as \(+1\) on the
translation coinvariants, which are \(R_r/(3)=\F_3\).

The norm operator of \(C_{3^t}\) is zero on (A.117.8):
it is a multiple of
\((1+x+\cdots+x^{d-1})s(x)\).
Thus \(H^1(C_{3^t},V)=V/(x-1)V=\F_3\).
The action on this cohomology includes
\(\epsilon^{-1}x\epsilon=x^{-1}\), so it is \(-1\).
The translation invariants of \(V\) also form one copy
of \(\F_3\) with sign \(-1\).
For clarity, on the free \(\Z_3\)-lattice
\(\Z_3[x]/s(x)\), multiplication by \(x-1\) has
cokernel \(\Z_3/3\); its Smith entries are units and a
single \(3\).  Modulo \(3^r\) its kernel has order three
and lies in \(3^{r-1}V\).
After division by \(3^{r-1}\), that line is
\((x-1)^{2d-1}\) in
\(\F_3[x-1]/(x-1)^{2d}\).
Both formulas (A.117.9) act on it by \(-1\).

The involution is central in the abelian cyclotomic group.
A group with such an involution acting as \(-1\) on a
three-primary module has zero cohomology in every degree,
by averaging over that order-two subgroup.
The two outside groups in the inflation--restriction
sequence therefore vanish, proving
\[
                         H^1(G_0,V)=0
                    \quad(V=C_r,\ C_r^\vee).
 \tag{A.117.10}
\]
This argument retains the larger cyclotomic field and the
deeper translation group.

\emph{The entire prior three-primary radical kernel occurs locally at three.}
Fix a completion \(R_w\) at three, and put
\(F_w=R_w(\theta)\).
The local group over \(\Q_3\) is
\(C_{3^t}\rtimes\operatorname{Gal}(R_w/\Q_3)\), with
the same literal splitting and cyclotomic action.
The element \(b\) is not a cube in \(F_w\).
Here is the exact quotient argument.
Already \(b\) is a noncube in \(\Q_3\), because its class
is that of \(-2\), which is \(7\) modulo nine; a
three-adic unit cube is \(1\) or \(-1\) modulo nine.
If \(b^{1/3}\) belonged to \(F_w\), its splitting field
\(\Q_3(b^{1/3},\mu_3)\) would be an \(S_3\) quotient
of this holomorph.
Every such quotient is the one belonging to
\(\Q_3(a^{1/3},\mu_3)\).
Indeed the translation must map onto the subgroup of
order three.  The image of the abelian cyclotomic group
contains a reflection, and is contained in its
centralizer, of order two.  Its image is therefore
determined by the cyclotomic sign; the kernel of the
translation map is its subgroup of index three.
These descriptions specify the same normal subgroup
for every surjection to \(S_3\).

Equality of those two splitting fields would make the
Kummer lines of \(a\) and \(b\) equal over \(\Q_3(\mu_3)\).
The valuation of \(a\) there is \(2\), whereas that
of \(b\) is zero.  The two nonzero multiples of the
first line have nonzero valuation modulo three.
The class of \(b\) remains nonzero after this quadratic
extension, since restriction on cubic classes is
injective by restriction followed by the degree-two norm.
This is a contradiction.

Now consider the actual map
\[
 \F_3[C_N]\longrightarrow F_w^\times/F_w^{\times3},
          \qquad (c_i)\longmapsto\prod_i\beta_i^{c_i}.
 \tag{A.117.11}
\]
Its kernel is a translation submodule.
Every nonzero submodule of
\(\F_3[C_N]=\F_3[Z]/(Z^N)\) contains its socle
\(Z^{N-1}=\sum_i x^i\).
The image of that socle is the actual product
\(\prod_i\beta_i=b\), just proved noncubic.
Thus (A.117.11) is injective.
Since the base contains all the specified roots of
unity, successive reduction of a power relation modulo
three proves independence at every needed three-power
precision.  In this reduction, the intervening roots
of unity are themselves cubes in the base up to that
precision.

The simultaneous \(\beta\)-root and \(b\)-root lattice
is exactly the lattice of A.105.7.
The independence just proved gives its full global
degree also over \(F_w\), including its common product
relation.  Adjoining the prime-to-three part of
\(b^{1/M}\) cannot change this: restriction on
three-power Kummer classes is injective under a
prime-to-three extension, by the norm map.
Consequently \(D_3\) contains the whole three-primary
factor of \(\operatorname{Gal}(E/F)\).
The remaining factor has order prime to three.
This statement does not require \(D_3\) to project
onto all of the global \(\operatorname{Gal}(R/\Q)\).

The kernel \(\operatorname{Gal}(E/F)\) acts trivially
on \(V\).  A cohomology class zero on \(D_3\) can
be represented by a cocycle zero there.
Its restriction to this kernel is a homomorphism.
It is zero on the entire three-primary factor by
the preceding paragraph, and on the prime-to-three
factor by the orders.  It therefore inflates from
\(H^1(G_0,V)=0\).  This proves (A.117.4).

\emph{The precise Poitou--Tate comparison.}
For a finite \(P\subset\mathcal P\), finite-module
Poitou--Tate duality
\cite[I, Theorem 4.10(b)]{BIR_MilneADT}
identifies the annihilator of the image on
\(S\setminus\{3\}\) with
\[
 \operatorname{res}_{S\setminus\{3\}}
 \left\{\xi\in H^1(G_{\Q,S\cup P},V^\vee):
       \operatorname{res}_3\xi=0,\
       \operatorname{res}_q\xi=0\ (q\in P)\right\}.
 \tag{A.117.12}
\]
The actual pairing is
\[
 ((c_w),\xi)\longmapsto
       \sum_{w\in S\setminus\{3\}}
           \operatorname{inv}_w(c_w\smile\operatorname{res}_w\xi).
 \tag{A.117.13}
\]
To obtain (A.117.12), extend the dual local tuple
by zero at three and at \(P\) and use orthogonality
of the two global images.
Every class inside the braces is unramified at \(P\),
so it lies in the fixed finite group
\(H^1(G_{\Q,S},V^\vee)\);
finiteness is
\cite[I, Corollary 4.15]{BIR_MilneADT}.
The decreasing kernels consequently stabilize.

Their intersection as \(P\) varies is
\[
 \ker\left(
 H^1(\operatorname{Gal}(E/\Q),V^\vee)
          \longrightarrow H^1(D_3,V^\vee)\right),
 \tag{A.117.14}
\]
viewed by inflation.
Indeed a global cocycle restricts on \(G_E\) to
a homomorphism, since \(E\) splits the module.
If it vanishes at every prime splitting in \(E\),
Chebotarev in its actual finite cocycle field makes
that homomorphism zero.  Inflation--restriction
then gives precisely the finite-group class in
(A.117.14).  The converse follows directly from
the splitting condition.  Removing any finite
set of test primes, or requiring a lower bound,
does not affect this Chebotarev argument.
By (A.117.4), (A.117.14) is zero.
The finite perfect local pairings therefore prove
(A.117.5), and stabilization proves that a finite
\(P\) suffices.

\emph{The global relative norm and the torus Kummer map.}
Set \(k=K_{n-1}\) and \(K=K_n\).
The Galois closure of \(K/k\) is \(K(\mu_3)\),
with group \(S_3\).
At the unique place of \(k\) over three, its
decomposition group is the whole \(S_3\): the
radicand \(\alpha^3\) has normalized valuation
one and the quadratic cyclotomic extension
is not contained in the odd-degree field \(k\).
The Hasse norm principle for this exact extension
follows from the full-decomposition-group argument
of A.77.  We recall its comparison to avoid replacing
the relative torus by an absolute norm torus.

Let \(T_0=R^1_{K/k}\mathbf G_m\), with character
lattice \(X\), and let \(L=K(\mu_3)\).
Every locally trivial class in \(H^2(k,X)\)
restricts over \(L\) to locally trivial characters,
using
\(H^2(G_L,\Z^m)=\operatorname{Hom}_{\rm cts}
 (G_L,\Q/\Z)^m\).
Chebotarev kills these characters.
Since \(H^1(G_L,X)=0\), the class inflates from
\(H^2(\operatorname{Gal}(L/k),X)\).
At the place with full decomposition group, local
inflation is injective by the identical \(H^1=0\)
argument, so that inflated class is zero.
Global torus duality consequently gives
\(\operatorname{Sha}^1(k,T_0)=0\);
see the torus duality paragraph in
\cite[I, \S4, p.~66]{BIR_MilneADT}, or the explicit
finite-group comparison in
\cite[Proposition 6.4 and Corollary 6.5]{BIR_Oki2026Norm}.
Now \(T_n=\operatorname{Res}_{k/\Q}T_0\).
Shapiro, together with its local product maps,
gives \(\operatorname{Sha}^1(\Q,T_n)=0\).

Apply (A.117.5) with \(V=C_r\) to the prescribed
local torus classes.  At three every class is
a local torus Kummer class: the relative local
cubic norm is onto by A.105, so
\(H^1(\Q_3,T_n)=0\).
At each \(q\in P\) the torus is split, and again
\(H^1(\Q_q,T_n)=0\).
At a place outside \(S\cup P\), the torus is
unramified and the produced class is unramified.
Such a class is a local Kummer class of an integral
torus point.  In fact reduction identifies its
unramified cohomology with finite-field cohomology;
the Kummer sequence there is onto in degree one
because \(H^1(\F_q,\mathcal T)=0\) for a connected
torus.  Hensel lifting lifts the resulting point
and its Kummer class.

Finally the exact global torus sequence is
\[
 0\longrightarrow T_n(\Q)/T_n(\Q)^{3^r}
 \xrightarrow{\ \kappa_{T_n,r}\ } H^1(\Q,C_r)
 \longrightarrow H^1(\Q,T_n)[3^r]\longrightarrow0.
 \tag{A.117.15}
\]
The constructed class maps to a locally trivial
torus torsor, hence to zero by the proved
Tate--Shafarevich vanishing.
It is therefore \(\kappa_{T_n,r}(U)\) for an
actual \(U\) with the literal norm (A.117.6).
This proves all the assertions.
\end{proof}

The unrestricted new local classes in this theorem
are still the pairs \((t_q,c_q)\) of A.109.
Its next-norm condition is the additional equation
\(c_q=z_qt_q+w_q\mathbf1\).
Neither the linear duality argument above nor the
absence of its finite dual obstruction proves that
equation.  In particular (A.117.5) does not prescribe
the three-adic class independently: its possible
values are the affine fibre cut out by the actual
pairings with
\(H^1(\operatorname{Gal}(E/\Q),C_r^\vee)\).

The retained certificate
\path{code/norm_prior_field_certificate.py}
checks the literal relative permutation modules,
their cyclic Smith invariants, finite cocycle
quotients for deeper holomorphs, and the unique
\(S_3\) quotient kernels.  Its output is retained
alongside the program.
